%% kapa: exam=2 round=36 year=2025 period=1 subject=appraisal_practice subject_ko=감정평가실무 seq=1 number=1 points=40
%% source: Q-Net 감정평가사 기출문제 (cst00309) artlSeq=5239325 file=raw/1008/5239325/2025년 제36회 감정평가사 2차 감정평가실무.pdf
%% conversion: kordoc 4.24.51 -> tools/gen2.py (auto); review=proofread-gpt-sol; flags=none
\begin{longquestion}{【문제 1】}
감정평가사 甲은 고객으로부터 건축공사가 중단된 부동산과 관련한 감정평가 및 관련 자문을 의뢰받았다. 관련 법규와 이론에 따라, 제시된 자료를 활용하여 다음 물음에 답하시오. (40점) \par
\medskip\noindent 물음 1) < 자료 1∼10 >을 활용하여 대상부동산에 대한 현재 상태대로의 적정 매수가격(2025.07.12. 시점)과 개발을 완료할 경우의 적정 매수가격(2026.07.12. 시점)을 제시하시오. (20점) \par
\medskip\noindent 물음 2) 개발을 완료할 경우의 적정 수익가격을 구하시오. (10점) \par
\medskip\noindent 물음 3) 주어진 자료에 따른 향후 시장동향을 기술하고, 이에 따른 수익성지수(PI), 내부수익률(IRR), 순현재가치(NPV)의 변화를 해당 구조식(산식)을 활용하여 설명하시오. (10점) \par
< 자료 1 > 기본적 사항 확정 \par
1. 기준가치: 시장가치 \par
2. 기준시점 \par
1) 2025.07.12. (현재시점) \par
2) 2026.07.12. (개발완료시점) \par
3. 대상부동산의 개요 \par
1) 토지 \par
{\small\centering
\resizebox{\ifdim\width>\linewidth\linewidth\else\width\fi}{!}{%
\begin{tabular}{|c|c|c|c|c|c|c|c|}
\hline
\makecell{소재지} & \makecell{지목} & \makecell{면적\\(m²)} & \makecell{용도\\지역} & \makecell{이용\\상황} & \makecell{도로\\접면} & \makecell{형상\\지세} & \makecell{주위환경} \\
\hline
\makecell{K시 H구\\A동 100} & \makecell{대} & \makecell{1,000} & \makecell{일반상업} & \makecell{업무용} & \makecell{광대한면} & \makecell{가장형\\평지} & \makecell{업무지대} \\
\hline
\end{tabular}}\par}

2) 건물: 지하3층, 지상7층으로 설계된 건물로, 현재 지하층 공사는 완료되었으나 지상층은 5층까지 주요구조부 공사 중 중단상태임 \par
{\small\centering
\resizebox{\ifdim\width>\linewidth\linewidth\else\width\fi}{!}{%
\begin{tabular}{|c|c|c|c|c|}
\hline
\makecell{소재지} & \makecell{구조} & \makecell{층} & \makecell{면적(m²)} & \makecell{용도} \\
\hline
\multirow{2}{*}{\makecell{K시 H구\\A동 100}} & \multirow{2}{*}{\makecell{철근콘크리트조\\슬래브지붕}} & \makecell{지상1∼7층} & \makecell{각층 800} & \makecell{업무용} \\
\cline{3-5}
 &  & \makecell{지하1∼3층} & \makecell{각층 800} & \makecell{주차장, 창고} \\
\hline
\end{tabular}}\par}

4. 표준지, 거래사례, 평가선례, 임대사례의 선정은 가장 비교성이 높은 것 하나만 선정하기로 함 \par
< 자료 2 > 공사중단 건물과 관련한 조사자료 \par
1. 기성공사 내역 \par
(단위: 천원) \par
{\small\centering
\resizebox{\ifdim\width>\linewidth\linewidth\else\width\fi}{!}{%
\begin{tabular}{|c|c|c|c|c|}
\hline
\multirow{2}{*}{\makecell{구분}} & \multicolumn{2}{|c|}{\makecell{설계안}} & \multicolumn{2}{|c|}{\makecell{기성공사}} \\
\cline{2-5}
 & \makecell{공사비} & \makecell{구성비(\%)} & \makecell{공사비} & \makecell{공정률(\%)} \\
\hline
\makecell{01. 가설공사} & \makecell{1,000,000} & \makecell{6.25} & \makecell{800,000} & \makecell{80} \\
\hline
\makecell{02. 기초 및 토공사} & \makecell{400,000} & \makecell{2.50} & \makecell{400,000} & \makecell{100} \\
\hline
\makecell{03. 철근콘크리트공사} & \makecell{2,800,000} & \makecell{17.50} & \makecell{2,240,000} & \makecell{80} \\
\hline
\makecell{04. 조적공사} & \makecell{200,000} & \makecell{1.25} & \makecell{160,000} & \makecell{80} \\
\hline
\makecell{05. 방수공사} & \makecell{400,000} & \makecell{2.50} & \makecell{320,000} & \makecell{80} \\
\hline
\makecell{06. 미장공사} & \makecell{600,000} & \makecell{3.75} & \makecell{180,000} & \makecell{30} \\
\hline
\makecell{07. 타일공사} & \makecell{600,000} & \makecell{3.75} & \makecell{180,000} & \makecell{30} \\
\hline
\makecell{08. 창호공사} & \makecell{1,800,000} & \makecell{11.25} & \makecell{540,000} & \makecell{30} \\
\hline
\makecell{09. 유리공사} & \makecell{400,000} & \makecell{2.50} & \makecell{120,000} & \makecell{30} \\
\hline
\makecell{10. 도장공사} & \makecell{2,000,000} & \makecell{12.50} & \makecell{600,000} & \makecell{30} \\
\hline
\makecell{11. 수장공사} & \makecell{1,600,000} & \makecell{10.00} & \makecell{480,000} & \makecell{30} \\
\hline
\makecell{12. 기타공사} & \makecell{600,000} & \makecell{3.75} & \makecell{480,000} & \makecell{80} \\
\hline
\makecell{소계} & \makecell{12,400,000} & \makecell{(77.50)} & \makecell{6,500,000} & \makecell{(52.41)} \\
\hline
\makecell{제경비} & \makecell{2,600,000} & \makecell{16.25} & \makecell{2,080,000} & \makecell{80} \\
\hline
\makecell{건축 공사비 합계} & \makecell{15,000,000} & \makecell{(93.75)} & \makecell{8,580,000} & \makecell{(57.20)} \\
\hline
\makecell{설계비} & \makecell{240,000} & \makecell{1.50} & \makecell{240,000} & \makecell{100} \\
\hline
\makecell{감리비} & \makecell{160,000} & \makecell{1.00} & \makecell{80,000} & \makecell{50} \\
\hline
\makecell{전기기본설비비} & \makecell{600,000} & \makecell{3.75} & \makecell{600,000} & \makecell{100} \\
\hline
\makecell{총 공사비 합계} & \makecell{16,000,000} & \makecell{100} & \makecell{9,500,000} & \makecell{(59.37)} \\
\hline
\end{tabular}}\par}

% [원문 그대로] 제조달원가
주) 동 공사비 내역서는 적정한 것으로 판단되어 2025년 7월 12일자 제조달원가로 적용할 수 있으며, 건축공사비와 공정률은 층별로 동일(기초공사 및 토공사는 공통 비용)한 것으로 가정함 \par
2. 공사중단에 따른 감가비용 등 관련 조사자료 \par
기성공사 중 공사중단에 따른 일부 물리적, 기능적 감가에 따른 비용은 기성 건축공사비의 5\% 정도로 추정되며 이는 적정한 것으로 조사되었고 건축허가사항과의 상이점 및 구조안전진단 결과 등에서 문제점이 발견되지는 않았으며, 이 경우 토지에 대한 별도 감가는 불필요함 \par
< 자료 3 > 향후 개발방안 관련 조사자료 \par
1. 기존 건축물 추가공사 방안 \par
1년이 소요되는 추가공사를 진행할 경우의 공사완료시점 재조달원가는 원자재비 및 금리 인상 등의 시장여건 변화를 반영하되, 설계안을 기준하여 미성 건축공사비의 20\% 증가요인으로 보정하는 것이 적정한 것으로 조사되었고, 정상적인 사용승인을 득하는데도 문제가 없을 것으로 조사됨 \par
2．기존 건축물 철거 후 신축방안 \par
매수인이 부담할 철거비용은 기성 건축공사비의 10\%가 적정한 수준이며, 신축시 신자재 및 신공법 도입으로 총 공사기간은 1년으로 동일하나, 총 공사비는 설계안 대비 25\% 증가요인이 있고 이는 적정한 것으로 조사되었고, 정상적인 사용승인을 득하는데도 문제가 없을 것으로 조사됨 \par
< 자료 4 > 표준지 공시지가 자료 \par
1. 표준지 공시지가 \par
(공시기준일: 2025.1.1.) \par
{\small\centering
\resizebox{\ifdim\width>\linewidth\linewidth\else\width\fi}{!}{%
\begin{tabular}{|c|c|c|c|c|c|c|c|c|}
\hline
\makecell{기호} & \makecell{소재지} & \makecell{지목} & \makecell{면적\\(m²)} & \makecell{이용\\상황} & \makecell{용도지역} & \makecell{도로접면} & \makecell{형상\\지세} & \makecell{공시지가\\(원/m²)} \\
\hline
\makecell{1} & \makecell{H구 A동\\88} & \makecell{대} & \makecell{1,000} & \makecell{업무용} & \makecell{근린상업} & \makecell{광대세각} & \makecell{정방형\\평지} & \makecell{6,600,000} \\
\hline
\makecell{2} & \multirow{2}{*}{\makecell{H구 A동\\102}} & \makecell{대} & \makecell{750} & \makecell{업무용} & \makecell{일반상업\\근린상업} & \makecell{광대한면} & \multirow{2}{*}{\makecell{부정형\\평지}} & \makecell{6,000,000} \\
\hline
\makecell{3} & \makecell{H구 A동\\147} & \makecell{대} & \makecell{900} & \makecell{상업용} & \makecell{일반상업} & \makecell{중로한면} & \makecell{정방형\\평지} & \makecell{4,500,000} \\
\hline
\makecell{4} & \makecell{H구 A동\\180} & \makecell{대} & \makecell{1,200} & \makecell{업무용} & \makecell{중심상업} & \makecell{광대소각} & \makecell{정방형\\평지} & \makecell{10,000,000} \\
\hline
\makecell{5} & \makecell{H구 A동\\196} & \makecell{대} & \makecell{900} & \makecell{업무용} & \makecell{일반상업} & \makecell{광대소각} & \makecell{정방형\\평지} & \makecell{7,000,000} \\
\hline
\end{tabular}}\par}

2. 표준지 조사자료 \par
기호 3 표준지는 위치를 확인한 바 후면 상업지대 토지였음 \par
< 자료 5 > 평가선례 자료(평가목적: 일반거래(시가참고용)) \par
{\small\centering
\resizebox{\ifdim\width>\linewidth\linewidth\else\width\fi}{!}{%
\begin{tabular}{|c|c|c|c|c|c|c|c|c|c|}
\hline
\makecell{기호} & \makecell{소재지} & \makecell{지목} & \makecell{면적\\(m²)} & \makecell{용도지역} & \makecell{이용\\상황} & \makecell{도로접면} & \makecell{형상\\지세} & \makecell{기준\\시점} & \makecell{토지단가\\(원/m²)} \\
\hline
\makecell{a} & \makecell{H구A동\\123} & \makecell{대} & \makecell{1,500} & \makecell{일반상업} & \makecell{업무용} & \makecell{광대한면} & \makecell{세장형\\평지} & \makecell{2024.01.01.} & \makecell{8,500,000} \\
\hline
\makecell{b} & \makecell{H구A동\\188} & \makecell{대} & \makecell{1,300} & \makecell{일반상업} & \makecell{상업용} & \makecell{광대한면} & \makecell{가장형\\평지} & \makecell{2023.01.01.} & \makecell{8,000,000} \\
\hline
\makecell{c} & \makecell{H구A동\\215} & \makecell{대} & \makecell{950} & \makecell{중심상업} & \makecell{업무용} & \makecell{광대한면} & \makecell{세장형\\평지} & \makecell{2021.01.01.} & \makecell{9,000,000} \\
\hline
\end{tabular}}\par}

< 자료 6 > 거래사례 자료 \par
1. 거래사례 \par
{\small\centering
\resizebox{\ifdim\width>\linewidth\linewidth\else\width\fi}{!}{%
\begin{tabular}{|c|c|c|c|c|c|c|c|c|c|}
\hline
\makecell{기호} & \makecell{소재지} & \makecell{지목} & \makecell{면적\\(m²)} & \makecell{용도지역} & \makecell{이용\\상황} & \makecell{도로접면} & \makecell{형상\\지세} & \makecell{거래\\시점} & \makecell{토지단가\\(원/m²)} \\
\hline
\makecell{가} & \makecell{H구A동\\87} & \makecell{대} & \makecell{290} & \makecell{일반상업} & \makecell{업무용} & \makecell{중로한면} & \makecell{정방형\\평지} & \makecell{2024.01.01.} & \makecell{5,550,000} \\
\hline
\makecell{나} & \multirow{2}{*}{\makecell{H구A동\\125}} & \makecell{대} & \makecell{3,000} & \makecell{일반상업\\3종일주} & \makecell{업무용} & \makecell{광대한면} & \multirow{2}{*}{\makecell{세장형\\평지}} & \makecell{2025.01.01.} & \makecell{8,500,000} \\
\hline
\makecell{다} & \makecell{H구A동\\190} & \makecell{대} & \makecell{2,000} & \makecell{일반상업} & \makecell{업무용} & \makecell{광대한면} & \makecell{가장형\\평지} & \makecell{2023.01.01.} & \makecell{8,000,000} \\
\hline
\makecell{라} & \makecell{H구A동\\220} & \makecell{대} & \makecell{950} & \makecell{일반상업} & \makecell{업무용} & \makecell{광대한면} & \makecell{세장형\\평지} & \makecell{2021.01.01.} & \makecell{9,000,000} \\
\hline
\makecell{마} & \makecell{H구A동\\363} & \makecell{대} & \makecell{1,160} & \makecell{일반상업} & \makecell{업무용} & \makecell{광대한면} & \makecell{가장형\\평지} & \makecell{2025.01.01.} & \makecell{8,500,000} \\
\hline
\end{tabular}}\par}

2. 거래사례 조사자료 \par
1) 기호 가 거래사례는 위치를 확인한 바 후면 상업지대 토지였음 \par
2) 기호 다 거래사례는 등기사항전부증명서를 확인한 바 미등재상태였음 \par
< 자료 7 > 시점수정 자료: 지가변동률(K시 H구) \par
(단위: \%) \par
{\small\centering
\resizebox{\ifdim\width>\linewidth\linewidth\else\width\fi}{!}{%
\begin{tabular}{|c|c|c|}
\hline
\makecell{구분} & \makecell{주거지역} & \makecell{상업지역} \\
\hline
\makecell{2021. 01. 01.～2025. 06. 30.(누계)} & \makecell{2.624} & \makecell{2.645} \\
\hline
\makecell{2022. 01. 01.～2025. 06. 30.(누계)} & \makecell{2.265} & \makecell{2.285} \\
\hline
\makecell{2023. 01. 01.～2025. 06. 30.(누계)} & \makecell{1.827} & \makecell{1.845} \\
\hline
\makecell{2024. 01. 01.～2025. 06. 30.(누계)} & \makecell{1.278} & \makecell{1.293} \\
\hline
\makecell{2025. 01. 01.～2025. 06. 30.(누계)} & \makecell{0.795} & \makecell{0.806} \\
\hline
\makecell{2025. 06. 01.～2025. 06. 30.} & \makecell{0.000} & \makecell{0.000} \\
\hline
\end{tabular}}\par}

주) 2025년 7월 이후 지가변동률은 미고시 되었음 \par
< 자료 8 > 토지 지역요인 \par
대상, 표준지, 평가선례 및 거래사례는 인근지역에 소재하여 지역요인은 대등함 \par
< 자료 9 > 토지 개별요인 \par
1. 가로조건(각지인 경우 가로조건에서 반영하기로 함) \par
{\small\centering
\resizebox{\ifdim\width>\linewidth\linewidth\else\width\fi}{!}{%
\begin{tabular}{|c|c|c|c|c|c|c|c|}
\hline
\makecell{구분} & \makecell{광대한면} & \makecell{광대소각} & \makecell{광대세각} & \makecell{중로한면} & \makecell{중로각지} & \makecell{소로한면} & \makecell{소로각지} \\
\hline
\makecell{광대한면} & \makecell{1.00} & \makecell{1.09} & \makecell{1.05} & \makecell{0.95} & \makecell{0.99} & \makecell{0.85} & \makecell{0.89} \\
\hline
\makecell{광대소각} & \makecell{0.92} & \makecell{1.00} & \makecell{0.96} & \makecell{0.87} & \makecell{0.91} & \makecell{0.78} & \makecell{0.82} \\
\hline
\makecell{광대세각} & \makecell{0.95} & \makecell{1.04} & \makecell{1.00} & \makecell{0.90} & \makecell{0.94} & \makecell{0.81} & \makecell{0.85} \\
\hline
\makecell{중로한면} & \makecell{1.05} & \makecell{1.15} & \makecell{1.11} & \makecell{1.00} & \makecell{1.04} & \makecell{0.89} & \makecell{0.94} \\
\hline
\makecell{중로각지} & \makecell{1.01} & \makecell{1.10} & \makecell{1.06} & \makecell{0.96} & \makecell{1.00} & \makecell{0.86} & \makecell{0.90} \\
\hline
\makecell{소로한면} & \makecell{1.18} & \makecell{1.28} & \makecell{1.24} & \makecell{1.12} & \makecell{1.16} & \makecell{1.00} & \makecell{1.05} \\
\hline
\makecell{소로각지} & \makecell{1.12} & \makecell{1.22} & \makecell{1.18} & \makecell{1.07} & \makecell{1.11} & \makecell{0.96} & \makecell{1.00} \\
\hline
\end{tabular}}\par}

2. 획지조건 \par
{\small\centering
\resizebox{\ifdim\width>\linewidth\linewidth\else\width\fi}{!}{%
\begin{tabular}{|c|c|c|c|c|c|}
\hline
\makecell{구분} & \makecell{정방형} & \makecell{가장형} & \makecell{세장형} & \makecell{사다리형} & \makecell{부정형} \\
\hline
\makecell{정방형} & \makecell{1.00} & \makecell{1.00} & \makecell{0.98} & \makecell{0.95} & \makecell{0.92} \\
\hline
\makecell{가장형} & \makecell{1.00} & \makecell{1.00} & \makecell{0.98} & \makecell{0.95} & \makecell{0.92} \\
\hline
\makecell{세장형} & \makecell{1.02} & \makecell{1.02} & \makecell{1.00} & \makecell{0.97} & \makecell{0.94} \\
\hline
\makecell{사다리형} & \makecell{1.05} & \makecell{1.05} & \makecell{1.03} & \makecell{1.00} & \makecell{0.97} \\
\hline
\makecell{부정형} & \makecell{1.09} & \makecell{1.09} & \makecell{1.07} & \makecell{1.03} & \makecell{1.00} \\
\hline
\end{tabular}}\par}

< 자료 10 > 기타 사항 \par
1. 제시된 이외의 조건은 동일한 것으로 가정함 \par
2. 지역요인비교치 및 개별요인비교치는 소수점 넷째자리에서 반올림하여 셋째자리까지 표시함 \par
3. 그 밖의 요인 보정치는 표준지 기준 산정방식을 적용하고 소수점 셋째자리 이하는 절사함 \par
4. 대상토지의 결정단가는 백원단위에서 반올림하여 천원단위까지 표시함 \par
< 자료 11 > 업무용빌딩 수익 관련자료 \par
1. 업무용빌딩 소득수익률(K시 H구) \par
{\small\centering
\resizebox{\ifdim\width>\linewidth\linewidth\else\width\fi}{!}{%
\begin{tabular}{|c|c|c|c|c|}
\hline
\makecell{구분} & \makecell{2024.3분기} & \makecell{2024.4분기} & \makecell{2025.1분기} & \makecell{2025.2분기} \\
\hline
\makecell{소득수익률(\%)} & \makecell{3.4} & \makecell{3.3} & \makecell{3.2} & \makecell{3.2} \\
\hline
\end{tabular}}\par}

주) 소득수익률 = $\dfrac{NOI}{V_{0}}$ ($NOI$: 순영업소득, $V_{0}$: 분기초 자산가액) \par
2. 업무용빌딩 영업경비율(K시 H구) \par
{\small\centering
\resizebox{\ifdim\width>\linewidth\linewidth\else\width\fi}{!}{%
\begin{tabular}{|c|c|c|c|c|}
\hline
\makecell{구분} & \makecell{2024.3분기} & \makecell{2024.4분기} & \makecell{2025.1분기} & \makecell{2025.2분기} \\
\hline
\makecell{영업경비율(\%)} & \makecell{48} & \makecell{49} & \makecell{50} & \makecell{50} \\
\hline
\end{tabular}}\par}

3. 임대사례 자료 \par
{\small\centering
\resizebox{\ifdim\width>\linewidth\linewidth\else\width\fi}{!}{%
\begin{tabular}{|c|c|c|c|c|c|}
\hline
\makecell{구분} & \makecell{대상물건} & \makecell{임대사례 1} & \makecell{임대사례 2} & \makecell{임대사례 3} & \makecell{임대사례 4} \\
\hline
\makecell{소재지} & \makecell{H구 A동} & \makecell{H구 A동} & \makecell{H구 A동} & \makecell{H구 B동} & \makecell{H구 B동} \\
\hline
\makecell{오피스 권역} & \makecell{KG서부\\권역} & \makecell{KG서부\\권역} & \makecell{KG서부\\권역} & \makecell{KG서부\\권역} & \makecell{KG동부\\권역} \\
\hline
\makecell{용도지역} & \makecell{일반상업} & \makecell{일반상업} & \makecell{일반상업} & \makecell{일반상업} & \makecell{일반상업} \\
\hline
\makecell{층(지상/지하)} & \makecell{7F/B3} & \makecell{18F/B3} & \makecell{9F/B3} & \makecell{8F/B3} & \makecell{25F/B6} \\
\hline
\makecell{구조} & \makecell{철근\\콘크리트} & \makecell{철골철근\\콘크리트} & \makecell{철근\\콘크리트} & \makecell{철근\\콘크리트} & \makecell{철골철근\\콘크리트} \\
\hline
\makecell{건물연면적(m²)} & \makecell{8,000} & \makecell{45,000} & \makecell{8,500} & \makecell{9,000} & \makecell{59,000} \\
\hline
\makecell{토지면적(m²)} & \makecell{1,000} & \makecell{3,300} & \makecell{1,260} & \makecell{1,200} & \makecell{3,200} \\
\hline
\makecell{사용승인(년)} & \makecell{2026.7.12.} & \makecell{2020.7.12.} & \makecell{2010.7.12.} & \makecell{2025.7.12.} & \makecell{2015.7.12.} \\
\hline
\makecell{연간실질임대료\\(원/m²)} & \makecell{-} & \makecell{221,000} & \makecell{185,000} & \makecell{213,000} & \makecell{252,000} \\
\hline
\makecell{임대(계약)시점} & \makecell{2026.7.12.} & \makecell{2025.7.12.} & \makecell{2025.7.12.} & \makecell{2025.7.12.} & \makecell{2025.7.12.} \\
\hline
\end{tabular}}\par}

주) 임대면적은 건물연면적으로 계약함 \par
4. 임대료지수 \par
{\small\centering
\resizebox{\ifdim\width>\linewidth\linewidth\else\width\fi}{!}{%
\begin{tabular}{|c|c|c|}
\hline
\makecell{구분} & \makecell{2025.7.12.} & \makecell{2026.7.12.} \\
\hline
\makecell{임대료지수} & \makecell{100} & \makecell{98} \\
\hline
\end{tabular}}\par}

5. 지역 및 개별요인 통합비교치 \par
{\small\centering
\resizebox{\ifdim\width>\linewidth\linewidth\else\width\fi}{!}{%
\begin{tabular}{|c|c|c|c|c|c|}
\hline
\multirow{2}{*}{\makecell{구분}} & \multirow{2}{*}{\makecell{대상}} & \multicolumn{4}{|c|}{\makecell{임대사례}} \\
\cline{3-6}
 &  & \makecell{1} & \makecell{2} & \makecell{3} & \makecell{4} \\
\hline
\makecell{평점} & \makecell{100} & \makecell{102} & \makecell{98} & \makecell{98} & \makecell{102} \\
\hline
\end{tabular}}\par}

6. 대상부동산의 순영업소득(NOI) 등과 관련하여 전문AI프로그램에 가격자료와 관련자료를 제시하고 계량분석을 실시한 결과 다음 결과를 얻었음 \par
1) 인근지역 업무시설의 최근 5년간 시계열 임대사례를 이용하여 대상부동산 순영업소득을 종속변수와 다수의 설명변수로 구축한 모형의 설명력은 0.306($R^{2}$ : 수정된 결정계수)이고, 모형과 설명변수들의 F값과 $t$값의 유의확률이 유의수준 보다 커서 추정된 모형이 유의하다는 가설을 기각하였음 \par
2) 인근지역 업무시설의 최근 5년간 시계열 임대사례를 이용하여 구축한 임대료변동률 추정모형의 설명력은 0.825( $R^{2}$ : 수정된 결정계수)이고, 모형과 독립변수의 F값과 $t$값의 유의확률이 유의수준 보다 작아 추정된 모형이 유의하다는 가설을 채택 하였으나 이를 장기간의 추이로 보기에는 어려움이 있다는 심사 감정평가사의 의견이 있었음 \par
$\hat{g}_{t}=-0.002g_{t-1}+5.125$ \par
($\hat{g}_{t}$: $t$시점의 임대료변동률 추정치, $g_{t-1}$: $t-1$기의 임대료변동률) \par
7. 자본환원표 \par
1) 이자율 연 4\% \par
{\small\centering
\resizebox{\ifdim\width>\linewidth\linewidth\else\width\fi}{!}{%
\begin{tabular}{|c|c|c|c|c|c|c|}
\hline
\makecell{기간(년)} & \makecell{일시불\\내가계수} & \makecell{연금\\내가계수} & \makecell{감채기금\\계수} & \makecell{일시불\\현가계수} & \makecell{연금\\현가계수} & \makecell{저당상수} \\
\hline
\makecell{1} & \makecell{1.040000} & \makecell{1.000000} & \makecell{1.000000} & \makecell{0.961538} & \makecell{0.961538} & \makecell{1.040000} \\
\hline
\makecell{5} & \makecell{1.216653} & \makecell{5.416323} & \makecell{0.184627} & \makecell{0.821927} & \makecell{4.451822} & \makecell{0.224627} \\
\hline
\end{tabular}}\par}

2) 이자율 연 5\% \par
{\small\centering
\resizebox{\ifdim\width>\linewidth\linewidth\else\width\fi}{!}{%
\begin{tabular}{|c|c|c|c|c|c|c|}
\hline
\makecell{기간(년)} & \makecell{일시불\\내가계수} & \makecell{연금\\내가계수} & \makecell{감채기금\\계수} & \makecell{일시불\\현가계수} & \makecell{연금\\현가계수} & \makecell{저당상수} \\
\hline
\makecell{1} & \makecell{1.050000} & \makecell{1.000000} & \makecell{1.000000} & \makecell{0.952381} & \makecell{0.952381} & \makecell{1.050000} \\
\hline
\makecell{5} & \makecell{1.276282} & \makecell{5.525631} & \makecell{0.180975} & \makecell{0.783526} & \makecell{4.329477} & \makecell{0.230975} \\
\hline
\end{tabular}}\par}

8. 기타 사항 \par
소득수익율과 영업경비율은 2025.2분기 이후 변동이 없다고 봄 \par
\end{longquestion}
